\documentclass[12pt]{article}

% Math packages
\usepackage{amsmath}
\usepackage{amssymb}
\usepackage{amsthm}
\usepackage{mathtools}

% Layout and formatting
\usepackage{geometry}
\geometry{margin=1in}
\usepackage{parskip}
\usepackage{microtype}

% Theorems and proofs
\theoremstyle{plain}
\newtheorem{theorem}{Theorem}
\newtheorem{lemma}[theorem]{Lemma}
\newtheorem{corollary}[theorem]{Corollary}
\theoremstyle{definition}
\newtheorem{definition}[theorem]{Definition}
\theoremstyle{remark}
\newtheorem{remark}[theorem]{Remark}

% Title information
\title{Exponential Separation of Singular Values in Random Matrix Products}
\author{}
\date{\today}

\begin{document}

\maketitle

\section{Proof: Exponential Separation of Singular Values in Random Matrix Products}

\subsection{Setup}
Let $X_1, X_2, ..., X_n$ be independent random matrices in $\mathbb{R}^{d \times d}$ with i.i.d. entries drawn from a distribution with mean 0 and variance $\sigma^2$.

Let $P_n = X_n X_{n-1} ... X_1$ be their product.

\subsection{Key Steps}

\subsubsection{Singular Value Decomposition}
For any matrix $A$, we can write its SVD as $A = U\Sigma V^T$ where:
\begin{itemize}
    \item $U, V$ are orthogonal matrices
    \item $\Sigma$ is diagonal with singular values $\sigma_1(A) \geq \sigma_2(A) \geq ... \geq \sigma_d(A) \geq 0$
\end{itemize}

\subsubsection{Lyapunov Exponents}
The key to understanding the asymptotic behavior is through Lyapunov exponents. For the product $P_n$, these are defined as:

\[
\lambda_i = \lim_{n \to \infty} \frac{1}{n} \log \sigma_i(P_n)
\]

\subsubsection{Oseledets Multiplicative Ergodic Theorem}
This theorem guarantees that these limits exist and are non-random under mild conditions. Moreover:

\begin{enumerate}
    \item The Lyapunov exponents are ordered: $\lambda_1 \geq \lambda_2 \geq ... \geq \lambda_d$
    \item With probability 1, at least $\lambda_1 - \lambda_d > 0$
\end{enumerate}

\subsubsection{Quantitative Bounds}
For i.i.d. random matrices with finite moments:

\[
\sigma_1(P_n) \approx e^{n\lambda_1}
\]
\[
\sigma_d(P_n) \approx e^{n\lambda_d}
\]

Therefore, the condition number grows exponentially:

\[
\kappa(P_n) = \frac{\sigma_1(P_n)}{\sigma_d(P_n)} \approx e^{n(\lambda_1 - \lambda_d)}
\]

\subsubsection{Explicit Values for Gaussian Case}
When entries are i.i.d. Gaussian with variance $\sigma^2/d$:

\[
\lambda_1 \approx \log(2) + \frac{1}{2}\log(\sigma^2)
\]
\[
\lambda_d \approx -\log(2) + \frac{1}{2}\log(\sigma^2)
\]

Thus $\lambda_1 - \lambda_d \approx 2\log(2)$, giving exponential growth of condition number:

\[
\kappa(P_n) \approx 4^n
\]

\subsection{Technical Details for Convergence}

The convergence relies on several key properties:

\begin{enumerate}
    \item \textbf{Furstenberg-Kesten theorem} states that for any norm $\|\cdot\|$:
    \[
    \lim_{n \to \infty} \frac{1}{n} \log \|P_n\| = \lambda_1 \text{ almost surely}
    \]

    \item The \textbf{multiplicative ergodic theorem} provides a filtration:
    \[
    \{0\} = V_d \subset V_{d-1} \subset ... \subset V_1 = \mathbb{R}^d
    \]
    where vectors in $V_k \setminus V_{k+1}$ grow at rate $\lambda_k$

    \item For any vector $v \in \mathbb{R}^d$, with probability 1:
    \[
    \lim_{n \to \infty} \frac{1}{n} \log \|P_n v\| = \lambda_k
    \]
    where $k = \max\{i : v \in V_i\}$
\end{enumerate}

\subsection{Implications}

This exponential separation has several important implications:

\begin{enumerate}
    \item The product becomes increasingly ill-conditioned
    \item Most directions in the input space get mapped to an exponentially small subspace
    \item Gradient-based optimization becomes challenging due to the exponential spread of singular values
\end{enumerate}

These results help explain why deep linear networks are difficult to train and motivate architectural choices like skip connections and normalization layers in modern neural networks.

\end{document}